\documentclass[a4paper,9pt]{extarticle}

%=========================================================
% PACKAGES
%=========================================================

\usepackage[
a4paper,
top=1.8cm,
bottom=1.5cm,
footskip=7mm,
left=1cm,
right=1cm,
headheight=14pt,
headsep=5mm
]{geometry}
\usepackage[T1]{fontenc}
\usepackage[utf8]{inputenc}
\usepackage[french]{babel}

\usepackage[sfdefault]{FiraSans}
\renewcommand{\familydefault}{\sfdefault}

\usepackage{multicol}
\setlength{\columnsep}{0.7cm}

\usepackage{amsmath,amssymb}

\usepackage{titlesec}
\usepackage{enumitem}

\usepackage{tcolorbox}
\tcbuselibrary{skins,breakable}

\usepackage{xcolor}

\usepackage{graphicx}
\usepackage{float}

\usepackage{tikz}

\usepackage{fancyhdr}
\usepackage[hidelinks]{hyperref}


%=========================================================
% PAGE SETTINGS
%=========================================================

\setlength{\parindent}{0pt}
\setlength{\parskip}{1.5pt}

\setlength{\abovedisplayskip}{4pt}
\setlength{\belowdisplayskip}{4pt}

\setlist[itemize]{
leftmargin=*,
itemsep=1pt,
topsep=1pt
}


%=========================================================
% COLORS
%=========================================================

\definecolor{Blue}{HTML}{E3F0FF}
\definecolor{Green}{HTML}{E3F7EA}
\definecolor{Pink}{HTML}{FFE5ED}
\definecolor{Orange}{HTML}{FFF0DC}
\definecolor{Purple}{HTML}{EEE7FF}
\definecolor{Yellow}{HTML}{FFFADB}
\definecolor{Teal}{HTML}{E0F5F3}


\definecolor{BlueTitle}{HTML}{285B8F}
\definecolor{GreenTitle}{HTML}{287052}
\definecolor{PinkTitle}{HTML}{A84063}
\definecolor{OrangeTitle}{HTML}{A96500}
\definecolor{PurpleTitle}{HTML}{6245A5}
\definecolor{TealTitle}{HTML}{1F6F68}


%=========================================================
% HEADER
%=========================================================

\pagestyle{fancy}

\fancyhf{}

\fancyhead[L]{\small\textbf{Modèles stochastiques}}
\fancyhead[C]{\small Vecteurs aléatoires et gaussiens}
\fancyhead[R]{\small Week 02 --- 14-20/09/2026}

\fancyfoot[C]{\small\href{https://github.com/AlphaXT38G/notes_ba5}{github.com/AlphaXT38G/notes\_ba5}}
\fancyfoot[R]{\small \thepage}

\renewcommand{\headrulewidth}{0.4pt}



%=========================================================
% TITLES
%=========================================================

\titleformat{\section}
{\large\bfseries\color{BlueTitle}}
{}{0pt}{}

\titlespacing{\section}
{0pt}{6pt}{3pt}


% separator

\newcommand{\separator}{
\vspace{2pt}
\hrule
\vspace{3pt}
}


%=========================================================
% BOXES
%=========================================================

% Boxes are kept whole (unbreakable): inside multicols, a box split by tcolorbox
% leaves an empty column behind it.

% DEFINITION - blue
\newtcolorbox{definition}{
enhanced,
unbreakable,
colback=Blue,
colframe=BlueTitle,
boxrule=0.7pt,
arc=0mm,
title=\textbf{Definition},
fonttitle=\bfseries
}

% IMPORTANT - pink
\newtcolorbox{important}{
enhanced,
unbreakable,
colback=Pink,
colframe=PinkTitle,
boxrule=0.7pt,
arc=0mm,
title=\textbf{Key idea},
fonttitle=\bfseries
}

% FORMULAS - green
\newtcolorbox{formula}{
enhanced,
unbreakable,
colback=Green,
colframe=GreenTitle,
boxrule=0.7pt,
arc=0mm,
title=\textbf{Key formulas},
fonttitle=\bfseries
}

% WARNING - orange
\newtcolorbox{warning}{
enhanced,
unbreakable,
colback=Orange,
colframe=OrangeTitle,
boxrule=0.7pt,
arc=0mm,
title=\textbf{Warning},
fonttitle=\bfseries
}

% EXAMPLE - purple
\newtcolorbox{example}{
enhanced,
unbreakable,
colback=Purple,
colframe=PurpleTitle,
boxrule=0.7pt,
arc=0mm,
title=\textbf{Example},
fonttitle=\bfseries
}

% SUMMARY - orange
\newtcolorbox{summary}{
enhanced,
unbreakable,
colback=Yellow,
colframe=OrangeTitle,
boxrule=0.7pt,
arc=0mm,
title=\textbf{Summary},
fonttitle=\bfseries
}



% SIMPLE - teal: the same idea again, in the plainest words, with an everyday picture
\newtcolorbox{simple}{
enhanced,
unbreakable,
colback=Teal,
colframe=TealTitle,
boxrule=0.7pt,
arc=0mm,
title=\textbf{Expliqué simplement},
fonttitle=\bfseries
}



%=========================================================
% FIGURE BOX
%=========================================================

\newtcolorbox{figurebox}{
enhanced,
colback=white,
colframe=gray,
boxrule=0.5pt,
arc=0mm,
center
}



%=========================================================
% DOCUMENT
%=========================================================

\begin{document}

%\footnotesize   --> To reduce the size of the text

\thispagestyle{fancy}



\begin{multicols}{2}


%=========================================================
\section{Où on en est}


\begin{summary}

\begin{itemize}

\item Mardi et vendredi : \textbf{module 2}, vecteurs aléatoires, puis vecteurs gaussiens.

\item Jeudi : séance de rattrapage (optionnelle) sur les notions de traitement du signal
(COM-202) utiles pour COM-300, donnée par Salim Najib. Son résumé est à la fin de ces notes.

\item Exercices : M1 8, 9 et 13 ; M2 1 et 4.

\end{itemize}

\end{summary}


\begin{simple}

La semaine passée : \textbf{un} nombre au hasard. Cette semaine : \textbf{plusieurs} nombres
au hasard en même temps (la taille et le poids d'une personne, le bruit sur deux antennes).
La question nouvelle : comment ils sont \textbf{liés} entre eux.

\end{simple}


%=========================================================
\section{Paire de v.a. : répartition et densité jointes}


\begin{formula}

\textbf{Fonction de répartition jointe} :

\[
F_{XY}(x,y) = P(X\le x, Y \le y)
\]

\textbf{Marginales} : $F_X(x) = \lim_{y\to\infty} F_{XY}(x,y)$, $F_Y(y) = \lim_{x\to\infty}
F_{XY}(x,y)$.

\vspace{3pt}

\textbf{Densité jointe} :

\[
f_{XY}(x,y) = \frac{\partial^2 F_{XY}(x,y)}{\partial x\,\partial y}, \qquad \iint f_{XY} = 1
\]

\[
P(a_1 < X \le a_2, b_1 < Y \le b_2) = \int_{a_1}^{a_2}\!\!\int_{b_1}^{b_2} f_{XY}(x,y)\,dy\,dx
\]

\textbf{Densités marginales} : on intègre l'autre variable.

\[
f_X(x) = \int_{-\infty}^\infty f_{XY}(x,y)\,dy \qquad f_Y(y) = \int_{-\infty}^\infty
f_{XY}(x,y)\,dx
\]

\end{formula}


\begin{formula}

\textbf{Cas discret}, $p_{ij} = P(X = x_i, Y = y_j)$ :

\[
f_{XY}(x,y) = \sum_{i,j} p_{ij}\,\delta(x-x_i)\,\delta(y-y_j)
\]

\end{formula}


\begin{definition}

$X$ et $Y$ sont \textbf{indépendantes} si $F_{XY}(x,y) = F_X(x)F_Y(y)$, d'où $f_{XY}(x,y) =
f_X(x)f_Y(y)$ (cas discret : $P(X = a, Y = b) = P(X=a)P(Y=b)$).

\vspace{3pt}

Si $X$ et $Y$ sont indépendantes, $g(X)$ et $h(Y)$ le sont aussi.

\end{definition}


\begin{simple}

La densité jointe, c'est le tas de sable posé \textbf{sur une table} (le plan $x,y$) au lieu
d'une ligne. La marginale de $X$, c'est ce qu'on voit en regardant le tas de côté : on écrase
tout le sable le long de $y$. Indépendantes : la forme du tas est juste « profil en $x$ fois
profil en $y$ », savoir $y$ ne dit rien sur $x$.

\end{simple}


%=========================================================
\section{Fonctions de deux v.a. et jacobien}


\begin{formula}

$Z = g_1(X,Y)$, $W = g_2(X,Y)$, et $(x_i,y_i)$ les racines de $z = g_1(x,y)$, $w = g_2(x,y)$.
Jacobien :

\[
J(x,y) = \det\begin{pmatrix} \frac{\partial g_1}{\partial x} & \frac{\partial g_1}{\partial
y} \\[2pt] \frac{\partial g_2}{\partial x} & \frac{\partial g_2}{\partial y}\end{pmatrix}
\]

\[
f_{ZW}(z,w) = \sum_i \frac{f_{XY}(x_i,y_i)}{|J(x_i,y_i)|}
\]

\end{formula}


\begin{formula}

\textbf{Cas linéaire} $\mathbf z = A\mathbf x + \mathbf b$, avec $\det A \ne 0$ : une seule
racine $\mathbf x = A^{-1}(\mathbf z - \mathbf b)$, et $J = \det A$ partout.

\[
f_{\mathbf Z}(\mathbf z) = \frac{f_{\mathbf X}(A^{-1}(\mathbf z - \mathbf b))}{|\det A|}
\]

\end{formula}


\begin{example}

\textbf{Somme de deux v.a. indépendantes} $Z = X+Y$. Il faut deux équations : on ajoute $W =
Y$ (qui ne nous intéresse pas), ce qui donne $A = \begin{pmatrix}1&1\\0&1\end{pmatrix}$, $\det
A = 1$. Donc $f_{ZW}(z,w) = f_{XY}(z-w, w) = f_X(z-w)f_Y(w)$, puis la marginale :

\[
f_Z(z) = \int_{-\infty}^\infty f_X(z-w)\,f_Y(w)\,dw \quad\Rightarrow\quad f_Z = f_X \star f_Y
\]

\end{example}


\begin{important}

La densité de la \textbf{somme} de deux v.a. \textbf{indépendantes} est le \textbf{produit de
convolution} de leurs densités.

\end{important}


\begin{simple}

Comme pour une seule v.a. : on change de coordonnées, et le jacobien dit de combien la
transformation \textbf{étire ou écrase les surfaces}. Si elle étire, le sable s'étale et sa
hauteur baisse d'autant. Astuce à retenir : quand on veut une seule nouvelle variable (comme
$X+Y$), on invente une deuxième bidon ($W = Y$) pour pouvoir utiliser la formule, puis on
l'intègre pour s'en débarrasser.

\end{simple}

\end{multicols}


\end{document}
